% GENERATED FILE. Edit canonical lessons, then run npm run generate:books.
% canonical-source-hash: fnv1a64:71890a453bbee828
% canonical-lessons: FA-C533-mojri, FA-C533-tanz, FA-C533-komedi, FA-C533-pakhsh-shodan, FA-C533-khire-shodan

\chapter{Shows and Presenters}
\label{ch:fa-shows-and-presenters}
\hlchaptermodality{533}

\begin{chapteropening}
\textbf{By the end of this chapter:} \emph{I can talk about TV shows and presenters.}

Complete the last lesson of chapter 533: I can talk about TV shows and presenters.
\end{chapteropening}

\section[\texorpdfstring{\fa{مجری} (mojri)}{مجری (mojri)}]{\fa{مجری} (mojri) — a presenter, a host (on TV or radio)}

\label{lesson:FA-C533-mojri}



\begin{quote}
Before the new one: say the Persian for a podcast, then the Persian for headphones.
\end{quote}

\subsection*{You'll want to know: \fa{مجری}}
\textbf{\fa{مجری}} (\emph{mojri}) — \textquotedblleft{}a presenter, a host (on TV or radio)\textquotedblright{}.

\begin{quote}
\fa{مجری} \fa{این} \fa{برنامه} \fa{خیلی} \fa{شوخ} \fa{است}.
\end{quote}

\begin{quote}
\emph{mojri-ye in barnâme kheyli shukh ast.}
\end{quote}

\begin{quote}
The presenter of this programme is very witty.
\end{quote}

Plural \emph{mojriyân} or \emph{mojrihâ}.

\begin{cousinweb}
\emph{mojri} is a loan, usually traced to Arabic \emph{mujrī}, 'one who carries something out'. EtymDB has no entry for it.
\end{cousinweb}

\begin{grammarlens}[title={What you've built}]
One more word for telling, explaining and describing at B1.
\end{grammarlens}

\subsection*{Your turn}
\begin{itemize}
\raggedright
  \item \emph{Say it:} \emph{mojri}
  \item \emph{Say it:} \emph{mojri}, once more
  \item \emph{Say it:} a sentence of your own with \emph{mojri}
  \item \emph{Recall:} say the Persian for a podcast, then the Persian for headphones, then say \emph{mojri} again
\end{itemize}

\begin{tcolorbox}[breakable,colback=teal!4,colframe=teal!35!black,title={Before you move on}]
What does \fa{مجری} mean? (\textquotedblleft{}A presenter, a host (on TV or radio)\textquotedblright{}.) Say it once more.
\end{tcolorbox}
\section[\texorpdfstring{\fa{طنز} (tanz)}{طنز (tanz)}]{\fa{طنز} (tanz) — satire, humour}

\label{lesson:FA-C533-tanz}



\begin{quote}
Before the new one: say the Persian for headphones, then the Persian for a presenter.
\end{quote}

\subsection*{You'll want to know: \fa{طنز}}
\textbf{\fa{طنز}} (\emph{tanz}) — \textquotedblleft{}satire, humour\textquotedblright{}.

\begin{quote}
\fa{این} \fa{نویسنده} \fa{داستانهای} \fa{طنز} \fa{مینویسد}.
\end{quote}

\begin{quote}
\emph{in nevisande dâstânhâ-ye tanz minevisad.}
\end{quote}

\begin{quote}
This writer writes humorous stories.
\end{quote}

\emph{tanz} — satire, comic writing. \emph{tanznevis} — a humorist.

\begin{cousinweb}
\emph{tanz} comes from Arabic \emph{ṭanz} (EtymDB), a word for mockery; Persian made it the name of satire.
\end{cousinweb}

\begin{grammarlens}[title={What you've built}]
One more word for telling, explaining and describing at B1.
\end{grammarlens}

\subsection*{Your turn}
\begin{itemize}
\raggedright
  \item \emph{Say it:} \emph{tanz}
  \item \emph{Say it:} \emph{tanz}, once more
  \item \emph{Say it:} a sentence of your own with \emph{tanz}
  \item \emph{Recall:} say the Persian for headphones, then the Persian for a presenter, then say \emph{tanz} again
\end{itemize}

\begin{tcolorbox}[breakable,colback=teal!4,colframe=teal!35!black,title={Before you move on}]
What does \fa{طنز} mean? (\textquotedblleft{}Satire, humour\textquotedblright{}.) Say it once more.
\end{tcolorbox}
\section[\texorpdfstring{\fa{کمدی} (komedi)}{کمدی (komedi)}]{\fa{کمدی} (komedi) — a comedy}

\label{lesson:FA-C533-komedi}



\begin{quote}
Before the new one: say the Persian for a presenter, then the Persian for satire.
\end{quote}

\subsection*{You'll want to know: \fa{کمدی}}
\textbf{\fa{کمدی}} (\emph{komedi}) — \textquotedblleft{}a comedy\textquotedblright{}.

\begin{quote}
\fa{ما} \fa{یک} \fa{کمدی} \fa{خنده} \fa{دار} \fa{دیدیم}.
\end{quote}

\begin{quote}
\emph{mâ yek komedi-ye khandedâr didim.}
\end{quote}

\begin{quote}
We watched a funny comedy.
\end{quote}

Plural \emph{komedihâ}.

\begin{cousinweb}
\emph{komedi} is a modern loan, usually traced to French \emph{comédie}. EtymDB has no entry for the Persian word.
\end{cousinweb}

\begin{grammarlens}[title={What you've built}]
One more word for telling, explaining and describing at B1.
\end{grammarlens}

\subsection*{Your turn}
\begin{itemize}
\raggedright
  \item \emph{Say it:} \emph{komedi}
  \item \emph{Say it:} \emph{komedi}, once more
  \item \emph{Say it:} a sentence of your own with \emph{komedi}
  \item \emph{Recall:} say the Persian for a presenter, then the Persian for satire, then say \emph{komedi} again
\end{itemize}

\begin{tcolorbox}[breakable,colback=teal!4,colframe=teal!35!black,title={Before you move on}]
What does \fa{کمدی} mean? (\textquotedblleft{}A comedy\textquotedblright{}.) Say it once more.
\end{tcolorbox}
\section[\texorpdfstring{\fa{پخش} \fa{شدن} (pakhsh shodan)}{پخش شدن (pakhsh shodan)}]{\fa{پخش} \fa{شدن} (pakhsh shodan) — to be broadcast; to spread}

\label{lesson:FA-C533-pakhsh-shodan}



\begin{quote}
Before the new one: say the Persian for satire, then the Persian for a comedy.
\end{quote}

\subsection*{You'll want to know: \fa{پخش} \fa{شدن}}
\textbf{\fa{پخش} \fa{شدن}} (\emph{pakhsh shodan}) — \textquotedblleft{}to be broadcast; to spread\textquotedblright{}.

\begin{quote}
\fa{این} \fa{سریال} \fa{هر} \fa{شب} \fa{از} \fa{تلویزیون} \fa{پخش} \fa{میشود}.
\end{quote}

\begin{quote}
\emph{in seriyâl har shab az televiziyon pakhsh mishavad.}
\end{quote}

\begin{quote}
This series is broadcast on TV every night.
\end{quote}

\emph{pakhsh} 'spreading' + \emph{shodan}. \emph{pakhsh kardan} — to broadcast; to hand out.

\begin{cousinweb}
\emph{pakhsh} is usually taken as a native Persian word for spreading, but EtymDB has no entry for it, and its origin is not settled. \emph{shodan} is inherited.
\end{cousinweb}

\begin{grammarlens}[title={What you've built}]
One more word for telling, explaining and describing at B1.
\end{grammarlens}

\subsection*{Your turn}
\begin{itemize}
\raggedright
  \item \emph{Say it:} \emph{pakhsh shodan}
  \item \emph{Say it:} \emph{pakhsh shodan}, once more
  \item \emph{Say it:} a sentence of your own with \emph{pakhsh shodan}
  \item \emph{Recall:} say the Persian for satire, then the Persian for a comedy, then say \emph{pakhsh shodan} again
\end{itemize}

\begin{tcolorbox}[breakable,colback=teal!4,colframe=teal!35!black,title={Before you move on}]
What does \fa{پخش} \fa{شدن} mean? (\textquotedblleft{}To be broadcast; to spread\textquotedblright{}.) Say it once more.
\end{tcolorbox}
\section[\texorpdfstring{\fa{خیره} \fa{شدن} (khire shodan)}{خیره شدن (khire shodan)}]{\fa{خیره} \fa{شدن} (khire shodan) — to stare}

\label{lesson:FA-C533-khire-shodan}



\begin{quote}
Before the new one: say the Persian for a comedy, then the Persian for to be broadcast.
\end{quote}

\subsection*{You'll want to know: \fa{خیره} \fa{شدن}}
\textbf{\fa{خیره} \fa{شدن}} (\emph{khire shodan}) — \textquotedblleft{}to stare\textquotedblright{}.

\begin{quote}
\fa{او} \fa{ساعتها} \fa{به} \fa{صفحه} \fa{نمایش} \fa{خیره} \fa{شد}.
\end{quote}

\begin{quote}
\emph{u sâ'athâ be safhe-ye namâyesh khire shod.}
\end{quote}

\begin{quote}
He stared at the screen for hours.
\end{quote}

\emph{be chizi khire shodan}; \emph{khire} 'dazzled, fixed' + \emph{shodan}.

\begin{cousinweb}
\emph{khire} means 'dazzled, fixed', but EtymDB has no entry for it, and its origin is not settled. \emph{shodan}, 'to become', is inherited.
\end{cousinweb}

\begin{grammarlens}[title={What you've built}]
One more word for telling, explaining and describing at B1.
\end{grammarlens}

\subsection*{Your turn}
\begin{itemize}
\raggedright
  \item \emph{Say it:} \emph{khire shodan}
  \item \emph{Say it:} \emph{khire shodan}, once more
  \item \emph{Say it:} a sentence of your own with \emph{khire shodan}
  \item \emph{Recall:} say the Persian for a comedy, then the Persian for to be broadcast, then say \emph{khire shodan} again
\end{itemize}

\begin{tcolorbox}[breakable,colback=teal!4,colframe=teal!35!black,title={Before you move on}]
What does \fa{خیره} \fa{شدن} mean? (\textquotedblleft{}To stare\textquotedblright{}.) Say it once more.
\end{tcolorbox}
